\section{Who Supplies the System's Values}
\label{sec:6.2}
A system's values have a supply chain. It begins before training, with
the conditions under which judgments are turned into labels and the
terms on which the underlying material is acquired.

Whose judgment is being encoded is a question about working conditions before it is a question about cognition. The labels that make a system safe are produced by people paid at the bottom of the supply chain to look at the material the system is being taught to refuse. To build the filter that keeps ChatGPT from reproducing the worst of its training data, OpenAI contracted with an outsourcing firm whose workers in Nairobi read and categorized descriptions of child sexual abuse, torture, bestiality, murder, and suicide, for take-home wages between roughly \$1.32 and \$2 an hour. The workers interviewed about it described lasting psychological injury; the counseling they were entitled to was rare and, by their account, useless against production quotas. The contract was terminated eight months early \autocite{perrigo2023openai}.

Two things follow. The first is the ordinary labor point: this is harmful work, priced as though it were not. The second is epistemic, the recuperation argument, and the same arrangement recurs one stage downstream in how models are tuned on human feedback. The point here is prior to both. A system's values are downstream of a labor arrangement, and the arrangement is a design choice.

The same question runs one stage further back. Labels are applied to material, and the material was collected. The industry's standard method was to take text and images at web scale without asking, and courts have begun to sort out what that means. A US ruling in Bartz v. Anthropic held that training a model on books the company had bought was transformative and protected as fair use, and that downloading the same books from piracy sites was not; the distinction turned on acquisition alone, and it produced the largest copyright settlement on record \autocite{bartz2025fairuse}. The defendant is named here on purpose: it is the company whose model helped produce this book, as the preface records. The European Union now requires the provider of a general-purpose model to publish a summary of what it trained on, against a template that asks for sources by category, scraped web content among them \autocite[Art.~53(1)(d)]{eu2024aiact}. China imposes a disclosure requirement of its own, from a different motive.

The argument that a web-scale corpus cannot be documented or audited at the size involved, and so encodes a hegemonic slice of the web while presenting as general, was made early and has not been answered by anything that came after \autocite{bender2021parrots}. A corpus assembled this way has no channel at all. An annotator has a queue and a guideline and could in principle be given a way to object; a writer or photographer whose work sits in a training set was never party to anything, so there is nothing for a channel to attach to. The annotation and feedback pipeline collects disagreement and moves none of it, and an obligation follows for a book making this argument. The stage described here sits before both, because no disagreement is collected where none is asked for.
